\documentclass[12pt]{report}
\usepackage{./sty/dissert}% style
\usepackage{comment}
\usepackage{hyperref}
\usepackage[acronym]{glossaries}
\usepackage{titlesec}
\usepackage{xcolor}
\usepackage{graphicx}
\usepackage{fancyhdr}
\usepackage{lastpage}
\usepackage{listings}
\usepackage{caption}
\usepackage{lmodern}
\usepackage[portuguese]{babel}
\usepackage{indentfirst}
\usepackage{float}
\usepackage{tabularx}

\lstdefinestyle{codigo}{
    basicstyle=\ttfamily\scriptsize,      
    columns=fullflexible,                 
    keepspaces=true,                      
    keywordstyle=\color{blue!70!black}\bfseries,
    commentstyle=\color{green!45!black}\itshape,
    stringstyle=\color{red!65!black},
    numbers=left, numberstyle=\tiny\color{gray}, numbersep=4pt,                        
    frame=single,
    framesep=2pt,                         
    aboveskip=4pt, belowskip=2pt,         
    xleftmargin=2pt, xrightmargin=2pt,    
    rulecolor=\color{gray!40},
    backgroundcolor=\color{gray!6},
    breaklines=true,
    showstringspaces=false,
    tabsize=2,                            
    captionpos=b,
    inputencoding=utf8,
    extendedchars=true,
    literate=
        {á}{{\'a}}1 {à}{{\`a}}1 {â}{{\^a}}1 {ã}{{\~a}}1
        {é}{{\'e}}1 {ê}{{\^e}}1
        {í}{{\'i}}1
        {ó}{{\'o}}1 {ô}{{\^o}}1 {õ}{{\~o}}1
        {ú}{{\'u}}1 {ü}{{\"u}}1
        {ç}{{\c{c}}}1
        {Á}{{\'A}}1 {À}{{\`A}}1 {Â}{{\^A}}1 {Ã}{{\~A}}1
        {É}{{\'E}}1 {Ê}{{\^E}}1
        {Í}{{\'I}}1
        {Ó}{{\'O}}1 {Ô}{{\^O}}1 {Õ}{{\~O}}1
        {Ú}{{\'U}}1
        {Ç}{{\c{C}}}1
        {º}{{\textordmasculine}}1 {ª}{{\textordfeminine}}1
        {°}{{\textdegree}}1 {€}{{\euro}}1
}
\lstset{style=codigo}
\captionsetup[lstlisting]{skip=2pt}

% Portuguese names
\renewcommand{\lstlistingname}{Código}
\renewcommand{\lstlistlistingname}{Lista de Códigos}
\hypersetup{
    colorlinks,
    linkcolor={black!50!black},
    citecolor={black!50!black},
    urlcolor={black!80!black},
    pdftitle={Ganhos Hands-On},
    pdfauthor={Nuno Costa PG63901, Fábio Alves PG63877},
    pdfdisplaydoctitle=true
}
\titleclass{\chapter}{straight}   % chapter no longer starts a new page

\titleformat{\chapter}{\LARGE}{\thechapter.}{20pt}{\LARGE\textbf}
\titlespacing*{\chapter}{0pt}{12pt}{10pt}   % space before / after
\renewcommand{\figurename}{Figura}
\titleformat*{\section}{\large\bfseries}
%
\makeglossaries%
 % ------------------- Document -----------------------------------------------------
\begin{document}
%\onehalfspacing
\pagenumbering{roman}
\newcommand\tab[1][1cm]{\hspace*{#1}}
%
\setlength{\headheight}{28pt}   
\pagestyle{fancy}
\fancyhf{}                   

% ---- Header ----
\fancyhead[L]{\includegraphics[height=0.8cm]{img/logo.png}}
\fancyhead[C]{Ganhos Hands-On}
\fancyhead[R]{CA - MEEIC}        

% ---- Footer (rodapé) ----
\fancyfoot[L]{Nuno Costa PG63901}
\fancyfoot[C]{Página \thepage\ de \pageref{LastPage}}
\fancyfoot[R]{Fábio Alves PG63877}

% ---- Lines ----
\renewcommand{\headrulewidth}{0.4pt}   
\renewcommand{\footrulewidth}{0.4pt}  
% --------------TRABALHO------------------
\pagenumbering{arabic}

% Ex 1 %%%%%%%%%%%%%%%%%%%%%%%%%%%%
\chapter{\texttt{aula2\_ganho\_L\_students.m}}

\textbf{TODO 1)} 

\begin{lstlisting}[language=Matlab, caption={$L_0 = p_d-p_c$}, label={lst:Lo}]
    Lo = pd(2:end).' - pc(2:end).';  
\end{lstlisting}

\textbf{TODO 2)} 

\begin{lstlisting}[language=Matlab, caption={$T = \mathcal{O}^{-1}\mathcal{O}_0$ e $L = TL_0$}, label={lst:T}]
    T = O^-1*Oo; 
    L1 = T*Lo; 
\end{lstlisting}

\textbf{TODO 3)} 

\begin{lstlisting}[language=Matlab, caption={[Fórmula de Ackermann]$L = p_d(A)\,\mathcal{O}^{-1}\,[0\ \cdots\ 0\ 1]^T$}, label={lst:ack}]
    L3 = polyvalm(pd,A)*O^-1*[0; 0; 1];
\end{lstlisting}

\textbf{TODO 4)} 

\begin{lstlisting}[language=Matlab, caption={Polos discretos $z=e^{pT_s}$}, label={lst:polos}]
    pz = exp(polos*Ts);
\end{lstlisting}

\textbf{Resultados:} 

\begin{lstlisting}[language=Matlab, caption={Resultados}, label={lst:results}, numbers=none]
    >> aula2_ganho_L_students
    eig(A) = [5.702 4.42e-16 -0.7016]
    -> a planta tem um polo em +5.70: e INSTAVEL.
    pd(s) coef.: [1 14 48 80]   pc(s) coef.: [1 -5 -4 1.768e-15]
    Lo = pd - pc = [19 52 80]
    M (Processo II) =
        0     1     1
        4     1    -4
        3     0    -3
            ProcI_canonica    ProcII_coef    ProcII_Ackermann     place      acker 
            ______________    ___________    ________________    _______    _______

        l1        100.33           100.33           100.33          100.33     100.33
        l2       -54.667          -54.667          -54.667         -54.667    -54.667
        l3        73.667           73.667           73.667          73.667     73.667
    eig(A-LC) = [-2+2i -2-2i -10]
    |e(4 s)|: malha aberta = 3.44e+08 ; observador = 0.002889
    Discreto (Ts=0.01 s): pz = [0.904837 0.980003+0.0196027i 0.980003-0.0196027i]
    Ld (coeficientes) = [0.91197 -0.48257 0.6694]
    Ld (place)        = [0.91197 -0.48257 0.6694]
    Ts*L (continuo)   = [1.0033 -0.54667 0.73667]   <- so uma aproximacao
    |eig(Ad-Ld*Cd)|   = [0.9802 0.9802 0.90484]
\end{lstlisting}

Os valores de $L_d$ e de $T_s L$ são próximos, mas não iguais, porque a discretização do observador não é exata. A aproximação é melhor quanto menor for $T_s$.

\clearpage

% Ex 2 %%%%%%%%%%%%%%%%%%%%%%%%%%%%
\chapter{\texttt{Oservadores\_aula2\_2627.m}}

\textbf{TODO 1)} 

\begin{lstlisting}[language=Matlab, caption={Ganho do controlador por feedback}, label={lst:ctrl}]
    K = acker(A,B,cfg.polosK);
\end{lstlisting}

\textbf{TODO 2)} 

\begin{lstlisting}[language=Matlab, caption={Ganho do observador}, label={lst:obs}]
    L = acker(A',C',cfg.polosL).';
\end{lstlisting}

\begin{table}[H]\centering
\caption{Efeito dos polos do observador (controlador em $\{-2\pm j,-2,-2\}$).}
\label{tab:polosL}
\begin{tabular}{ccccc}\hline
Polos $L$ & $L^T$ & $|e(8\,\text{s})|$ & $|x(8\,\text{s})|$ & $\max|u|$ [N$\cdot$m]\\\hline
$-1$  & $[4\;\;6.05\;\;0.43\;\;1.11]$        & $0.105$             & $1.08$               & $3.01$\\
$-2$  & $[8\;\;24.049\;\;3.302\;\;2.7318]$   & $0.000168$          & $0.0548$             & $2.82$\\
$-5$  & $[20\;\;150.05\;\;51.068\;\;65.441]$ & $4.14\times10^{-14}$& $0.000389$           & $14.8$\\
$-20$ & $[80\;\;2400\;\;3262\;\;16323]$      & $2.1\times10^{-17}$ & $2.0\times10^{-5}$   & $369$\\\hline
\end{tabular}
\end{table}

\textbf{Observadores:}
\begin{table}[H]\centering
\caption{Efeito dos polos do observador (controlador em $\{-2\pm j,-2,-2\}$), em relação ao caso de referência ($-5$).}
\label{tab:polosL}
\small
\begin{tabularx}{\textwidth}{p{2.7cm}XX}\hline
Comparação & Observador lento ($-1$) & Observador rápido ($-20$)\\\hline

Erro de estimação &
Converge devagar; ainda $|e(8\,\text{s})|=0.105$. &
Anula-se quase de imediato ($|e(8\,\text{s})|\approx10^{-17}$). \\[2pt]

Posição máxima da bola &
Maior: o controlador atua sobre uma estimativa errada e a bola ainda não estabilizou aos 8~s ($|x|=1.08$). &
Semelhante à de $-5$; resposta ditada pelo controlador ($|x(8\,\text{s})|\approx2\times10^{-5}$). \\[2pt]

Ângulo máximo da viga &
Maior e mais prolongado, para corrigir o desvio da bola. &
Picos bruscos no início, causados pelos picos de binário. \\[2pt]

Binário máximo &
Baixo: $\max|u|=3.01$~N$\cdot$m. &
Muito elevado: $\max|u|=369$~N$\cdot$m (saturaria o atuador e amplifica o ruído). \\\hline
\end{tabularx}
\end{table}

\textbf{Observador lento ($-1$ vs $-2$):} como os polos da malha fechada são os do controlador mais os do observador. 
Com $-1$, os polos do observador são os mais lentos e dominam a resposta: o controlador atua durante muito tempo sobre uma estimativa errada.
A bola afasta-se mais do centro, oscila e demora a estabilizar: aos 8~s ainda $|e|=0.105$ e $|x|=1.08$.

\textbf{Observador rápido ($-20$):} o erro de estimação anula-se quase de imediato, mas o preço é um ganho $L$ enorme ($[80;\,2400;\,3262;\,16323]$).
O erro inicial é amplificado, o que gera picos de binário de $369$~N$\cdot$m (contra $3$~N$\cdot$m com $-1$): o atuador saturaria e qualquer ruído na medida de $x$ seria fortemente amplificado. 

\textbf{Animação:} com $-1$, a bola estimada (vermelho tracejado) demora a
alcançar a real (verde), e a viga inclina-se com base nessa estimativa
errada, por isso a bola real percorre mais a viga. Com $-20$, as duas
coincidem quase desde o início, mas a viga faz um movimento brusco nos
primeiros instantes, consequência dos picos de binário.

%-----------------------------------

\if false   
\begin{lstlisting}[language=Matlab, caption={Resultados}, label={lst:result}, numbers=none]
    >> Observadores_aula2_2627(struct('polosL', [-1 -1 -1 -1]))
    Modelo linearizado em x=0, u=0:
    A =
            0    1.0000         0         0
            0         0    9.8100         0
            0         0         0    1.0000
        0.9807         0    0.0490         0
    B =
            0
            0
            0
    -0.4998
    eig(A) = [-1.768 1.768 1.11e-16+1.754i 1.11e-16-1.754i]
    K = [-6.0409 -7.342 -50.116 -16.006]
    L = [4 6.049 0.42774 1.1128]
    |e(8 s)| = 0.105 ; |x(8 s)| = 1.08 ; max|u| = 3.01 N.m
    >> Observadores_aula2_2627(struct('polosL', [-20 -20 -20 -20]))
    Modelo linearizado em x=0, u=0:
    A =
            0    1.0000         0         0
            0         0    9.8100         0
            0         0         0    1.0000
        0.9807         0    0.0490         0
    B =
            0
            0
            0
    -0.4998
    eig(A) = [-1.768 1.768 1.11e-16+1.754i 1.11e-16-1.754i]
    K = [-6.0409 -7.342 -50.116 -16.006]
    L = [80 2400 3262.4 16323]
    |e(8 s)| = 2.09e-17 ; |x(8 s)| = 1.98e-05 ; max|u| = 369 N.m
    >> Observadores_aula2_2627(struct('polosL', [-5 -5 -5 -5]))
    Modelo linearizado em x=0, u=0:
    A =
            0    1.0000         0         0
            0         0    9.8100         0
            0         0         0    1.0000
        0.9807         0    0.0490         0
    B =
            0
            0
            0
    -0.4998
    eig(A) = [-1.768 1.768 1.11e-16+1.754i 1.11e-16-1.754i]
    K = [-6.0409 -7.342 -50.116 -16.006]
    L = [20 150.05 51.068 65.441]
    |e(8 s)| = 4.14e-14 ; |x(8 s)| = 0.000389 ; max|u| = 14.8 N.m
\end{lstlisting}
\fi






















\if false
\noindent
\begin{minipage}[t]{0.40\textwidth}
Para verificar a observabilidade calcula-se a matriz
$\mathcal{O}$ com \lstinline{obsv()} e compara-se a sua
característica com a ordem do sistema $n$. Como
$\mathrm{rank}(\mathcal{O}) = 2 = n$, o sistema é
observável.
\begin{lstlisting}[language=Matlab, caption={Resultado}, label={lst:obs_result}, numbers=none]
O =
     1     0
     0     1
r =
     2
O sistema é observável
\end{lstlisting}
\end{minipage}
\hfill
\begin{minipage}[t]{0.45\textwidth}
\begin{lstlisting}[language=Matlab]
O = obsv(A, C)
r = rank(O)
if r == n
    fprintf('O sistema é observável\n')
else
    fprintf('O sistema não é observável\n');
end
\end{lstlisting}
\captionof{lstlisting}{Observabilidade}
\label{lst:obsv}

\end{minipage}

\textbf{1.3)} Com $L=0$ a planta é estável ($\operatorname{eig}(A)=-0.75\pm4.41j$), logo o erro decai sozinho, sendo oscilatório, ao ritmo imposto pela planta.

\begin{figure}[htbp]
    \centering
    \begin{subfigure}[b]{0.48\textwidth}
        \centering
        \includegraphics[width=\linewidth]{img/logo.png}
        \label{fig:estados1}
    \end{subfigure}
    \hfill
    \begin{subfigure}[b]{0.48\textwidth}
        \centering
        \includegraphics[width=\linewidth]{img/logo.png}
        \label{fig:medicao1}
    \end{subfigure}
    \caption{Resultados da simulação com L = 0}
    \label{fig:Simul 1}
\end{figure}

% Ex 2 %%%%%%%%%%%%%%%%%%%%%%%%%%%%
\noindent \textbf{2)}

\textbf{2.1)} Em discreto, a estabilidade é definida pelo módulo dos polos ($|z|<1 =>$ Estável), porém -3 e -4 são polos contínuos, e aplicando a regra: $|-3|>1$ e $|-4|>1$, logo os polos estão fora do círculo unitário, indicando instabilidade.

\textbf{2.2)}
\noindent
\begin{minipage}[t]{0.40\textwidth}
Com os pólos escolhidos, a matriz de ganho do observador é calculada com \lstinline{place()}.
\begin{lstlisting}[language=Matlab, caption={Resultado}, label={lst:obs_result}, numbers=none]
pd =
    0.9704    0.9608
L =
    0.0519
   -0.1663
\end{lstlisting}
\end{minipage}
\hfill
\begin{minipage}[t]{0.45\textwidth}
\begin{lstlisting}[language=Matlab]
pd = exp(cfg.poles*cfg.Ts)
L = place(Ad',Cd',pd)'
\end{lstlisting}
\captionof{lstlisting}{Pólos e ganho do observador}
\label{lst:obsv}

\end{minipage}

\textbf{2.3)}
\begin{minipage}[t]{0.40\textwidth}
$eig$ é igual a $p_d$ (a menos da ordem e de erros numéricos $\approx10^{-10}$, pelo que a comparação se faz com \texttt{sort} e tolerância, não com \texttt{==}).
\begin{lstlisting}[language=Matlab, caption={Resultado}, label={lst:obs_result}, numbers=none]
ei =
    0.9608 
    0.9704
Os pólos coincidem
\end{lstlisting}
\end{minipage}
\hfill
\begin{minipage}[t]{0.45\textwidth}
\begin{lstlisting}[language=Matlab]
ei = eig(Ad-L*Cd)
if abs(sort(ei') - sort(pd)) < 1e-10
    fprintf('Os pólos coincidem\n')
else
    fprintf('Os pólos não coincidem\n')
end
\end{lstlisting}
\captionof{lstlisting}{Comparação de polos}
\label{lst:polos}
\end{minipage}

% Ex 3 %%%%%%%%%%%%%%%%%%%%%%%%%%%%
\noindent \textbf{3)}

\textbf{3.1--3.2)} 

\begin{lstlisting}[language=Matlab, caption={Inovação e atualização do observador}, label={lst:inov}]
innovation(j)=Y(j)-Cd*H(:,j);
H(:,j+1)=Ad*H(:,j)+Bd*u(j)+L*innovation(j);
\end{lstlisting}

\textbf{3.3)} O observador converge em $\approx1.5$~s, sem oscilações, enquanto a
malha aberta segue a dinâmica lenta da planta (Figura~\ref{fig:ex3}). A
velocidade é corrigida sem ser medida porque está acoplada à posição: um erro
em $\hat v$ desvia $\hat p$ da medida, o que aparece na inovação e é corrigido
por $L$.
\begin{figure}[htbp]
    \centering
    \begin{subfigure}[b]{0.45\textwidth}
        \centering
        \includegraphics[width=\linewidth]{img/logo.png}
        \label{fig:estados3}
    \end{subfigure}
    \hfill
    \begin{subfigure}[b]{0.45\textwidth}
        \centering
        \includegraphics[width=\linewidth]{img/logo.png}
        \label{fig:medicao3}
    \end{subfigure}
    \caption{Resultados da simulação com observador}
    \label{fig:ex3}
\end{figure}

\textbf{3.4)} $X(:,j)$ e $H(:,j)$ referem-se ao mesmo instante $t_j$.
        Comparar $X(:,j)$ com $H(:,j{+}1)$ introduziria um erro de
        desalinhamento $\approx T_s \dot{x}$, que não se anula (o sistema
        é sempre excitado por $u$), mesmo com um observador perfeito.

\textbf{Verificação:} após 10~s, $\|e\| = 5.07\times10^{-14} < 10^{-8}$. A convergência não implica $\hat{x}\to0$, pois o sistema continua a receber a entrada sinusoidal.

\noindent\textbf{4)}

\textbf{4.1)} Polos mais rápidos reduzem o transiente, mas exigem um ganho $L$ maior, amplificando o ruído, sobretudo na
velocidade.

\begin{table}[htbp]
    \centering
    \begin{minipage}[t]{0.35\textwidth}
        \hspace*{1em}\textbf{4.2)} O caso c converge mais depressa, mas tem pior RMSE em regime: é o compromisso rapidez--ruído.
    \end{minipage}
    \hfill
    \begin{minipage}[t]{0.60\textwidth}
        \vspace{0pt}
        \centering
        \caption{RMSE após 2 s.}
        \label{tab:polos}
        \begin{tabular}{lccc}\hline
        Caso & Polos & RMSE $p$ [m] & RMSE $v$ [m/s]\\\hline
        a ($\sigma=0$)    & $-3,-4$   & 6.181e-05 & 0.00016021\\
        b ($\sigma=0.03$) & $-3,-4$   & 0.0044406 & 0.023322\\
        c ($\sigma=0.03$) & $-12,-16$ & 0.011647  & 0.053313\\\hline
        \end{tabular}
    \end{minipage}
\end{table}

\begin{figure}[htbp]
    \centering
    \begin{minipage}[t]{0.45\textwidth}
        \vspace{0pt}
        \raggedright
        \hspace*{1em}\textbf{4.3)} \textbf{Azul:} massa real (só disponível em simulação); \textbf{verde:} estimativa
do observador; \textbf{vermelho:} réplica em malha aberta. Com ruído,
$y_j-C_d\hat x_j=C_de_j+n_j$: mesmo com $e=0$ a inovação fica igual ao ruído,
logo não tende para zero. 
    \end{minipage}
    \hfill
    \begin{minipage}[t]{0.45\textwidth}
        \vspace{0pt}
        \centering
        \includegraphics[width=\linewidth]{img/logo.png}
        \caption{Frame da animação.}
        \label{fig:estados}
    \end{minipage}
\end{figure}

\begin{figure}[htbp]
    \centering
    \begin{minipage}[t]{0.45\textwidth}
        \vspace{0pt}
        \raggedright
        \textbf{5)} Com $k_\text{real}=24$, subtraindo planta e observador: \[e_{j+1}=(A_d-LC_d)e_j+(A_r-A_d)x_j+(B_r-B_d)u_j\]
O termo forçado não se anula, porque $x$ e $u$ são sinusoidais, logo o erro
fica limitado mas não nulo (RMSE $\approx0.018$~m e $0.10$~m/s). Polos mais
rápidos reduzem-no, à custa de maior sensibilidade ao ruído.
    \end{minipage}
    \hfill
    \begin{minipage}[t]{0.45\textwidth}
        \vspace{0pt}
        \centering
        \includegraphics[width=\linewidth]{img/logo.png}
        \caption{Medição e inovação com $k_\text{real}=24$}
        \label{fig:estados}
    \end{minipage}
\end{figure}

\fi












\printglossaries


\end{document}
